\DocumentMetadata{
  pdfversion=2.0,
  pdfstandard=ua-2,
  lang=en-US,
  tagging=on,
  tagging-setup={math/setup=mathml-SE,tabsorder=structure}
}
\documentclass[11pt,letterpaper]{article}
\usepackage[margin=0.68in,includefoot]{geometry}
\usepackage{newcomputermodern}
\usepackage{amsmath,amscd,mathtools}
\usepackage{tikz,adjustbox}
\providecommand{\square}{\mdlgwhtsquare}
\usepackage[hidelinks]{hyperref}
\hypersetup{
  pdftitle={Analysis First-Year Exam, January 2025, Part 2},
  pdfauthor={University of Florida Department of Mathematics},
  pdfsubject={Graduate mathematics examination},
  pdfdisplaydoctitle=true,
  bookmarksopen=true
}
\setcounter{secnumdepth}{0}
\setlength{\parindent}{0pt}
\setlength{\parskip}{0.28em}
\setlength{\emergencystretch}{3em}
\setlength{\leftmargini}{2.15em}
\setlength{\leftmarginii}{2.65em}
\setlength{\leftmarginiii}{3em}
\pagestyle{empty}
\raggedbottom
\allowdisplaybreaks
\NewTaggingSocketPlug{block/list/label}{examproblem}{%
  \tagstructbegin{tag=Lbl}\tagstructend%
  \tagstructbegin{tag=\LBody}%
  \tagstructbegin{tag=\ProblemHeadingTag,title-o={\ProblemHeadingText}}%
  \tagmcbegin{tag=\ProblemHeadingTag,actualtext={\ProblemHeadingText}}%
  #2%
  \tagmcend%
  \tagstructend%
}
\newenvironment{examproblems}[2]{%
  \def\ProblemHeadingTag{#2}%
  \AssignTaggingSocketPlug{block/list/label}{examproblem}%
  \begin{enumerate}%
  \setcounter{enumi}{#1}%
  \renewcommand{\labelenumi}{\textbf{\theenumi.}}%
  \setlength{\topsep}{0.35em}%
  \setlength{\partopsep}{0pt}%
  \setlength{\itemsep}{0.48em}%
  \setlength{\parsep}{0.18em}%
}{\end{enumerate}}
\newenvironment{examsubparts}{%
  \AssignTaggingSocketPlug{block/list/label}{default}%
  \begin{enumerate}%
  \setlength{\topsep}{0.18em}%
  \setlength{\partopsep}{0pt}%
  \setlength{\itemsep}{0.16em}%
  \setlength{\parsep}{0.1em}%
}{\end{enumerate}}
\newenvironment{examinstructions}{%
  \par\begingroup\itshape
}{%
  \par\endgroup\vspace{0.18em}
}
\makeatletter
\renewcommand\subsection{\@startsection{subsection}{2}{0pt}%
  {-1.25ex plus -.25ex minus -.1ex}%
  {0.55ex plus .1ex}%
  {\normalfont\large\bfseries}}
\renewcommand{\@maketitle}{%
  \newpage\null\vskip -1em%
  {\footnotesize\bfseries
    \href{https://www.ufl.edu/}{University of Florida}, \href{https://math.ufl.edu/}{Department of Mathematics}\par}%
  \vskip -0.5em%
  \tagstructbegin{tag=H1,title={Analysis First-Year Exam, January 2025, Part 2}}%
  \tagpdfparaOff%
  \tagmcbegin{tag=H1}%
    {\LARGE\bfseries\href{https://gma.math.ufl.edu/exams/analysis-first-year/}{Analysis First-Year Exam}, January 2025, Part 2\par}%
  \tagmcend%
  \tagpdfparaOn%
  \tagstructend%
  \vskip 0.25em%
  \tagmcbegin{artifact}%
    \hrule height 1pt%
  \tagmcend%
  \vskip 1em%
}
\makeatother
\title{Analysis First-Year Exam, January 2025, Part 2}
\author{}
\date{}
\begin{document}
\maketitle
\thispagestyle{empty}
\begin{examinstructions}
Answer FOUR questions in detail. State carefully any results used without proof.
\end{examinstructions}
\begin{examproblems}{0}{H2}
\def\ProblemHeadingText{Problem 1.}
\item
Let \(f:[0,1]\to\mathbb{R}\) be Riemann integrable. If~\(f\) is continuous at~\(0\), does it follow that
\[
\int_0^1 f(x^n)\,dx \to f(0) \quad\text{as } n\to\infty?
\]
 If so, give a proof; if no, give a counterexample.
\def\ProblemHeadingText{Problem 2.}
\item
Let \((f_n)_{n=0}^{\infty}\) be a sequence of continuous real-valued functions on the metric space~\(X\), in which \((x_n)_{n=0}^{\infty}\) is a convergent sequence.
\begin{examsubparts}
\item[\textnormal{(i)}]
Prove that if \((f_n)_{n=0}^{\infty}\) is uniformly convergent then \((f_n(x_n))_{n=0}^{\infty}\) converges.
\item[\textnormal{(ii)}]
Show by example that \((f_n(x_n))_{n=0}^{\infty}\) can fail to converge if \((f_n)_{n=0}^{\infty}\) is merely pointwise convergent.
\end{examsubparts}
\def\ProblemHeadingText{Problem 3.}
\item
Let~\(\Phi\) be an equicontinuous family of real-valued functions on the connected metric space~\(X\). Prove that if~\(\Phi\) is bounded at one point of~\(X\) then~\(\Phi\) is bounded at every point of~\(X\). [Note: to say that~\(\Phi\) is bounded at \(p\in X\) is to say that \(\{f(p):f\in\Phi\}\subseteq\mathbb{R}\) is bounded.]
\def\ProblemHeadingText{Problem 4.}
\item
Let \((f_n)_{n=0}^{\infty}\) be a sequence of measurable real-valued functions on some measure space. Prove that each of the following sets is measurable.
\begin{examsubparts}
\item[\textnormal{(i)}]
\(A=\{\omega:f_n(\omega)\to 0\text{ as }n\to\infty\}\).
\item[\textnormal{(ii)}]
\(B=\{\omega:\sum_n f_n(\omega)\text{ converges absolutely}\}\).
\end{examsubparts}
\def\ProblemHeadingText{Problem 5.}
\item
This problem has two parts.
\begin{examsubparts}
\item[\textnormal{(i)}]
State the Monotone Convergence Theorem.
\item[\textnormal{(ii)}]
On a measure space, let \((f_n)_{n=0}^{\infty}\) be a sequence of non-negative measurable functions converging pointwise to the function~\(f\). Prove that if there exists a real number~\(K\) such that each of the integrals \(\int f_n\,d\mu\) is at most~\(K\), then~\(f\) is integrable and \(\int f\,d\mu\leq K\).
\end{examsubparts}
\end{examproblems}
\end{document}
